%! Boxplots and Comparing Distributions — Unit 1, Lesson 1.6 — Guided Notes & Practice
% THE in-class document, in the MAIN IDEAS / NOTES shape (see LESSON_SHAPE.md §1): the
% vocabulary box, then ONE two-column guidednotes table. Each row is one idea — a label on the
% left; on the right a terse definition the student completes, then a grid of short numbered
% problems with reserved answer space. Problems are numbered 1..26 through the whole component.
%   I DO   : the teacher fills the definition lines and works the FIRST problem of each grid.
%   WE DO  : the class works the rest of each grid, students holding the pen; the last row,
%            Guided Practice, is worked entirely together in a fresh context.
%   YOU DO : nothing on this page — ../homework/ IS the individual practice, started in the
%            last minutes of the period. No independent set, no reflection box, no objectives
%            box (the targets are on the cover), no notesbox, no practicebox.
% REGENERATED FROM THE RETIRED EFFL SHAPE (LESSON_SHAPE.md §7): the old QuickNotes became these
% rows' definition fills and \stepnum procedure, the old Activity's Twenty Free Throws and Two
% Bus Routes became instruction rows 1-4, and the old Application: Two Science Sections became
% the Guided Practice row. The old Check Your Understanding items stay in ../homework/.
% VOCABULARY IS GIVEN HERE, NOT DISCOVERED — there is no spoiler rule any more.
%
% THE CRUX is row 3, EVERY SECTION HOLDS ONE QUARTER. Target misconception: "a wide section
% means most of the data is there." Problem 14 is the trap (the classmate who reads width as
% count) and problem 15 is the counterweight (width reports spread, never count).
%
% THE PROTECTED DATA (rows 1-3: free throws made out of 100, twenty players):
%   30 31 31 32 33 | 36 36 37 38 38 | 41 43 45 47 49 | 52 56 60 66 74
%   min 30 · Q1 34.5 · median 39.5 · Q3 50.5 · max 74
%   Section widths: 4.5, 5.0, 11.0, 23.5 — the right section is over five times the left —
%   yet each holds exactly FIVE of the twenty. That count is the crux (EK UNC-1.L.2).
%   Every boundary falls inside a gap in the data (33|36, 38|41, 49|52), so the counting in
%   problem 13 is unambiguous: 5, 5, 5, 5. PROTECTED — do not soften the width contrast.
% THE PROTECTED DATA (row 4, travel minutes): Route A 18/21/23/26/31 (IQR 5, no outliers);
%   Route B 16/20/25/32/58 (IQR 12; upper cut-off 32 + 1.5(12) = 50, so 58 is set apart and
%   the whisker stops at 40, the most extreme value that is NOT unusually far).
% THE PROTECTED DATA (Guided Practice, test points): Period 2 58/74/80/86/94 (IQR 12, skewed
%   left); Period 5 60/68/72/84/98 (IQR 16, skewed right). Neither has an outlier.
% CONTEXTS ARE DISTINCT BY DESIGN: the notes teach on free throws and bus routes, Guided
% Practice is the science sections, ../ap_practice/ is weekday-vs-weekend tips, and
% ../homework/ is reading minutes, bus delays, and pizza deliveries.
%
% PAGE LOCKSTEP with notes_key/: every \blank{} here becomes a short \ans{} there, every
% \termblank becomes a \termans, and every \pcell{n}{...}{H}{} gets its answer in the fourth
% argument. Heights and blank widths are sized to the answers so the two files set identically.
% The \bxp macro is authored BYTE-IDENTICALLY in the blank and the key, as is every TikZ block.
% TABLE MECHANICS: inside a cell, `\\` ENDS THE ROW — break lines with \par. A row cannot break
% across pages: each long idea is split at a seam (after the figure, before the prompt) with a
% bare `\\` and continued with `& ...` as a sub-row.
% NO SKETCH QUESTIONS: every display is pre-drawn in TikZ and read; students fill tables only.
\documentclass[12pt]{article}
\usepackage{apstats-article}
\usepackage{apstats-boxes}
% \bxp{y-center}{half-height}{min}{Q1}{median}{Q3}{max} — every boxplot on this page is
% PRE-DRAWN with it. Authored byte-identically in the blank and the key.
\newcommand{\bxp}[7]{%
  \draw[thick, navy] (#3,#1) -- (#4,#1);
  \draw[thick, navy] (#6,#1) -- (#7,#1);
  \draw[thick, navy] (#3,{#1-#2}) -- (#3,{#1+#2});
  \draw[thick, navy] (#7,{#1-#2}) -- (#7,{#1+#2});
  \draw[thick, navy, fill=skymid] (#4,{#1-#2}) rectangle (#6,{#1+#2});
  \draw[very thick, navy] (#5,{#1-#2}) -- (#5,{#1+#2});}

\begin{document}

\pageheader{Unit 1, Lesson 1.6}{Guided Notes \& Practice}
% NAMESTRIP: no \namedateperiod here — the name row lives on the cover only.

\vspace{2pt}

\boxguard[16]
\begin{vocabbox}
\small
Fill in each term \textbf{as we name it} in the notes below.
\par
\termblank{Five-number summary}
\par
\termblank{Boxplot}
\par
\termblank{Whisker}
\par
\termblank{Modified boxplot}
\par
\termblank{Parallel boxplots}
\par
\end{vocabbox}

\small
\begin{guidednotes}
% ── ROW 1: the five-number summary ───────────────────────────────────────────
\mainidea[The five]{Numbers} &
The \textbf{five-number summary} of a quantitative variable is, in order, the \blank{1.6cm},
\blank{0.9cm}, the \blank{1.5cm}, \blank{0.9cm}, and the \blank{1.6cm}. A \textbf{boxplot} is
that summary \blank{1.2cm} to scale on a number line --- five numbers, nothing else.

\vspace{3pt}
Twenty players on a squad; each made this many free throws out of 100 attempts. Every player is
one dot, and the boxplot below was built from those same twenty numbers.

\vspace{2pt}
\begin{center}
\begin{tikzpicture}[scale=0.8]
  \foreach \x in {0,1,2,3,4,5,6,7,8,9} {
    \draw[linegray, very thin] (\x,0.02) -- (\x,0.55);
    \draw[linegray, very thin] (\x,-1.70) -- (\x,-0.85);}
  \foreach \x/\n in {0/1, 0.2/2, 0.4/1, 0.6/1, 1.2/2, 1.4/1, 1.6/2, 2.2/1, 2.6/1, 3.0/1,
                     3.4/1, 3.8/1, 4.4/1, 5.2/1, 6.0/1, 7.2/1, 8.8/1} {
    \foreach \k in {1,...,\n} { \fill[navy] (\x,{0.2*\k}) circle (0.075); }}
  \draw[->, thick] (-0.35,0) -- (9.4,0);
  \foreach \x/\lab in {0/30, 1/35, 2/40, 3/45, 4/50, 5/55, 6/60, 7/65, 8/70, 9/75} {
    \draw (\x,-0.08) -- (\x,0.08); \node[below, font=\scriptsize] at (\x,-0.12) {\lab};}
  \bxp{-1.25}{0.36}{0}{0.9}{1.9}{4.1}{8.8}
  \node[font=\scriptsize] at (4.5,-1.98) {Free throws made (out of 100)};
\end{tikzpicture}
\end{center}
\\
& \notesprompt{Read the five numbers off the dotplot, then find each one in the boxplot.}
\begin{probgrid}{|Y|Y|Y|}
\pcell{1}{Minimum}{0.9cm}{} &
\pcell{2}{$Q_1$}{0.9cm}{} &
\pcell{3}{Median}{0.9cm}{} \\ \hline
\pcell{4}{$Q_3$}{0.9cm}{} &
\pcell{5}{Maximum}{0.9cm}{} &
\pcell{6}{Which part of the boxplot marks $Q_1$?}{1.3cm}{} \\ \hline
\end{probgrid}
\\ \hline
% ── ROW 2: boxplot anatomy, as a procedure ───────────────────────────────────
\mainidea[Anatomy of a]{Boxplot} &
The \textbf{box} runs from \blank{0.9cm} to \blank{0.9cm} --- the middle \blank{0.8cm}\% of the
data --- and the line inside it is the \blank{1.5cm}. Each \textbf{whisker} reaches the most
extreme value that is \blank{0.9cm} an outlier. An outlier is not swallowed by a whisker: it
gets its \blank{1.5cm}, which makes the display a \textbf{modified boxplot} --- the version the
AP exam draws.

\vspace{3pt}
\stepnum{1}\ Mark the five numbers on a \blank{2.2cm}.\par
\stepnum{2}\ Draw the box from $Q_1$ to $Q_3$, with a line at the \blank{1.5cm}.\par
\stepnum{3}\ Test each end with the \blank{2.0cm} rule from Lesson~1.5.\par
\stepnum{4}\ Run each whisker out to the most extreme value that is \blank{2.6cm}; plot any
outlier on its own.
\notesprompt{Use the free-throw boxplot above.}
\begin{probgrid}{|Y|Y|Y|}
\pcell{7}{How wide is the box?}{1.2cm}{} &
\pcell{8}{How many of the twenty players sit inside the box?}{1.2cm}{} &
\pcell{9}{Does this plot show an outlier? How can you tell?}{1.2cm}{} \\ \hline
\pcell{10}{Between which two totals does the top quarter sit?}{1.5cm}{} &
\pcell{11}{The lower whisker runs 30 to 34.5. What does its short length say?}{1.5cm}{} &
\pcell{12}{Name one thing the dotplot shows that the boxplot cannot.}{1.5cm}{} \\ \hline
\end{probgrid}
\\ \hline
% ── ROW 3: THE CRUX — every section holds one quarter ────────────────────────
\mainidea[Every section holds]{One quarter} &
The five numbers cut a boxplot into \blank{1.0cm} sections: lower whisker, lower half of the
box, upper half of the box, upper whisker. Each one holds \blank{0.8cm}\% of the data, however
wide it looks. Width reports \blank{1.4cm} --- it never reports \blank{1.2cm}.

\vspace{4pt}
\textbf{13.} Count the twenty dots in each section of the free-throw plot above. The widths are
given; every boundary falls in a gap, so each count is exact.

\vspace{2pt}
\footnotesize\setlength{\tabcolsep}{2pt}%
\renewcommand{\arraystretch}{1.25}
\begin{tabularx}{\linewidth}{@{} l Y Y Y Y @{}}
\rowcolor{sky}
\textbf{Section} & 30 to 34.5 & 34.5 to 39.5 & 39.5 to 50.5 & 50.5 to 74 \\
\hline
\textbf{Width (shots)} & 4.5 & 5.0 & 11.0 & 23.5 \\
\textbf{Players} & \blank{0.8cm} & \blank{0.8cm} & \blank{0.8cm} & \blank{0.8cm} \\
\end{tabularx}\small
\\
& \fcolorbox{navy}{sky}{\parbox{\dimexpr\linewidth-2\fboxsep-2\fboxrule\relax}{\centering
\textbf{The widest section is over five times the narrowest ---}\\
\textbf{and holds exactly the \blank{1.1cm} number of players.}}}
\notesprompt{Answer the classmate.}
\begin{probgrid}{|Y|Y|}
\pcell{14}{A classmate says, \emph{``the wide part is where most of the shots landed.''} Use
your counts to say what is wrong.}{1.5cm}{} &
\pcell{15}{So what does that wide section actually tell you about those players?}{1.5cm}{} \\ \hline
\pcell{16}{Ten players made how many free throws or fewer, and which of the five numbers told
you?}{1.5cm}{} &
\pcell{17}{The dotplot has a tight clump from 30 to 38. Why can the boxplot not show
it?}{1.5cm}{} \\ \hline
\end{probgrid}
\\ \hline
% ── ROW 4: comparing two distributions on a shared axis ─────────────────────
\mainidea[Comparing two groups]{Parallel Boxplots} &
Two or more boxplots drawn on one shared axis are \textbf{parallel} (side-by-side) boxplots.
Comparing distributions means addressing \blank{1.4cm}, \blank{2.0cm}, and unusual features ---
all three, every time, each with a \blank{2.2cm} word (\emph{higher than}, \emph{more variable
than}) and each \blank{1.9cm}. Describing one group and then describing the other is
\emph{not} a comparison and earns nothing.

\vspace{3pt}
Under skew the mean leaves the median: skewed right pulls the mean to the \blank{1.1cm} of the
median; skewed left pulls it to the \blank{0.9cm}.
\\
& A commuter timed the same morning trip on two bus routes.

\vspace{2pt}
\begin{center}
\begin{tikzpicture}[scale=0.85]
  \foreach \x in {0,0.8,1.6,2.4,3.2,4.0,4.8,5.6,6.4,7.2} {
    \draw[linegray, very thin] (\x,-0.10) -- (\x,1.80);}
  \bxp{1.40}{0.30}{0.48}{0.96}{1.28}{1.76}{2.56}
  \bxp{0.50}{0.30}{0.16}{0.80}{1.60}{2.72}{4.00}
  \fill[navy] (6.88,0.50) circle (0.085);
  \node[font=\bfseries\scriptsize, anchor=east] at (-0.22,1.40) {Route A};
  \node[font=\bfseries\scriptsize, anchor=east] at (-0.22,0.50) {Route B};
  \draw[->, thick] (-0.20,-0.10) -- (7.75,-0.10);
  \foreach \x/\lab in {0/15, 0.8/20, 1.6/25, 2.4/30, 3.2/35, 4.0/40, 4.8/45, 5.6/50, 6.4/55,
                       7.2/60} {
    \draw (\x,-0.18) -- (\x,-0.02); \node[below, font=\scriptsize] at (\x,-0.22) {\lab};}
  \node[font=\scriptsize] at (3.6,-0.95) {Travel time (minutes)};
\end{tikzpicture}
\end{center}

\footnotesize\setlength{\tabcolsep}{2pt}%
\renewcommand{\arraystretch}{1.2}
\begin{tabularx}{\linewidth}{@{} l Y Y Y Y Y @{}}
\rowcolor{sky}
\textbf{Route} & \textbf{Min} & \textbf{$Q_1$} & \textbf{Median} & \textbf{$Q_3$} &
\textbf{Max} \\
\hline
\textbf{A} & 18 & 21 & 23 & 26 & 31 \\
\textbf{B} & 16 & 20 & 25 & 32 & 58 \\
\end{tabularx}\small
\\
& \notesprompt{Compare the two routes.}
\begin{probgrid}{|Y|Y|}
\pcell{18}{How wide is the middle half of each route's times?}{1.4cm}{} &
\pcell{19}{Which route is typically faster, and which is less predictable? Say it in
context.}{1.4cm}{} \\ \hline
\pcell{20}{Route B's 58-minute morning is drawn on its own. Show why, and say where its whisker
stops.}{1.7cm}{} &
\pcell{21}{Route B is skewed right. Is its mean above or below its median of 25 minutes?
Why?}{1.7cm}{} \\ \hline
\end{probgrid}
\\ \hline
% ── LAST ROW: GUIDED PRACTICE — we do, together, a fresh context ─────────────
\mainidea[Guided practice]{Two Science Sections} &
Two sections of the same course sat one test. The five-number summaries are given, in points.

\vspace{2pt}
\begin{center}
\begin{tikzpicture}[scale=0.85]
  \foreach \x in {0,1.4,2.8,4.2,5.6,7.0} {
    \draw[linegray, very thin] (\x,-0.15) -- (\x,1.75);}
  \bxp{1.35}{0.30}{1.12}{3.36}{4.20}{5.04}{6.16}
  \bxp{0.50}{0.30}{1.40}{2.52}{3.08}{4.76}{6.72}
  \node[font=\bfseries\scriptsize, anchor=east] at (-0.22,1.35) {Period 2};
  \node[font=\bfseries\scriptsize, anchor=east] at (-0.22,0.50) {Period 5};
  \draw[->, thick] (-0.20,-0.15) -- (7.55,-0.15);
  \foreach \x/\lab in {0/50, 1.4/60, 2.8/70, 4.2/80, 5.6/90, 7.0/100} {
    \draw (\x,-0.23) -- (\x,-0.07); \node[below, font=\scriptsize] at (\x,-0.27) {\lab};}
  \node[font=\scriptsize] at (3.5,-1.00) {Test score (points)};
\end{tikzpicture}
\end{center}

\footnotesize\setlength{\tabcolsep}{2pt}%
\renewcommand{\arraystretch}{1.2}
\begin{tabularx}{\linewidth}{@{} l Y Y Y Y Y @{}}
\rowcolor{sky}
\textbf{Section} & \textbf{Min} & \textbf{$Q_1$} & \textbf{Median} & \textbf{$Q_3$} &
\textbf{Max} \\
\hline
\textbf{Period 2} & 58 & 74 & 80 & 86 & 94 \\
\textbf{Period 5} & 60 & 68 & 72 & 84 & 98 \\
\end{tabularx}\small
\\
& \notesprompt{We work these together.}
\begin{probgrid}{|Y|Y|}
\pcell{22}{Report each section's median, with units.}{1.2cm}{} &
\pcell{23}{Report each section's IQR, with the subtraction.}{1.2cm}{} \\ \hline
\pcell{24}{Which section is skewed, and which way? How does the picture tell you?}{1.7cm}{} &
\pcell{25}{Does either section show an outlier? How can you tell from the display?}{1.7cm}{} \\ \hline
\end{probgrid}
\\
& \begin{probgrid}{|Y|}
\pcell{26}{\textbf{One paragraph, in context}, comparing the two sections: center, variability,
and unusual features --- every claim carrying a comparative word.}{2.0cm}{} \\ \hline
\end{probgrid}
\\ \hline
\end{guidednotes}

% ── THE NOTES END AT GUIDED PRACTICE ─────────────────────────────────────────
% There is deliberately no "you do" here: ../homework/ is the individual practice,
% started in class in the last 8 minutes while the teacher circulates.

\end{document}
